Take $A=(0,0)$, $B=(100,0)$, $C=(100,100)$, $D=(0,100)$, and $M=(50,0)$. The two circles have equations $x^2+y^2-100x=0$ and $x^2+y^2-200y=0$. Subtraction gives $x=2y$, and substitution gives $y=0$ or $40$. The intersection other than $A$ is $E=(80,40)$. Line $CE$ has equation $y=3x-200$, so its intersection with $AB$ is $F=(\dfrac{200}{3},0)$. Thus $AF=\dfrac{200}{3}$ and $\boxed{m+n=203}$.
\begin{asy}
texpreamble("\usepackage{amsmath,amssymb}");
size(240);
defaultpen(fontsize(9pt));
pair A=(0,0),B=(100,0),C=(100,100),D=(0,100),M=(50,0),E=(80,40),F=(200/3.0,0);draw(A--B--C--D--cycle);draw(circle(M,50),blue);draw(circle(D,100),gray);draw(C--F,red);label("A",A,SW);label("B",B,SE);label("C",C,NE);label("D",D,NW);label("M",M,S);label("E",E,(1,0));label("F",F,S);
\end{asy}
